\chapter{Object-Oriented Programming in Python: Classes and Objects}

\section{Introduction: Real-Life Hook}

Imagine your school has multiple students. Each student has:

\begin{itemize}
    \item Name  
    \item Age  
    \item Grade  
\end{itemize}

Instead of creating separate lists for each attribute, \textbf{OOP lets us create a ``Student'' blueprint}. Each student becomes an \textbf{object} of this blueprint (class), making programs \textbf{organized and reusable}.

\section{Concept: What is a Class and an Object?}

\begin{definition}{Class}
A \textbf{class} is a blueprint for creating objects. It defines attributes (data) and methods (functions) that each object will have.
\end{definition}

\begin{definition}{Object}
An \textbf{object} is an instance of a class. Each object has the attributes and methods defined by its class.
\end{definition}

\section{Creating a Class}

\begin{verbatim}
class Student:
    def __init__(self, name, age, grade):
        self.name = name
        self.age = age
        self.grade = grade

    def greet(self):
        print(f"Hello, my name is {self.name} and I am {self.age} years old.")
\end{verbatim}

- \texttt{\_\_init\_\_} is the \textbf{constructor}, automatically called when creating an object  
- \texttt{self} refers to the object itself  

\section{Creating Objects}

\begin{verbatim}
# Creating student objects
student1 = Student("Amina", 15, "9th")
student2 = Student("Rafi", 14, "8th")

# Access attributes
print(student1.name)  # Amina
print(student2.grade) # 8th

# Call methods
student1.greet()
student2.greet()
\end{verbatim}

\section{Modifying Object Attributes}

\begin{verbatim}
student1.grade = "10th"
print(student1.grade)  # 10th
\end{verbatim}

- Object attributes can be updated individually  

\section{Mini-Project: Student Database}

Create a program to store multiple students and print their details:

\begin{verbatim}
class Student:
    def __init__(self, name, age, grade):
        self.name = name
        self.age = age
        self.grade = grade

    def display(self):
        print(f"{self.name}, Age: {self.age}, Grade: {self.grade}")

students = [
    Student("Amina", 15, "9th"),
    Student("Rafi", 14, "8th"),
    Student("Sadia", 16, "10th")
]

for student in students:
    student.display()
\end{verbatim}

Output:

\begin{verbatim}
Amina, Age: 15, Grade: 9th
Rafi, Age: 14, Grade: 8th
Sadia, Age: 16, Grade: 10th
\end{verbatim}

\section{Step-by-Step Diagram: Class and Object Flow}

\begin{center}
\begin{tikzpicture}[node distance=1.5cm]
\node (class) [draw, rectangle, fill=yellow!20] {Define Class (Student)};
\node (object1) [draw, ellipse, fill=green!20, below left of=class, xshift=-2.5cm, yshift=-1cm] {Object student1};
\node (object2) [draw, ellipse, fill=green!20, below right of=class, xshift=2.5cm, yshift=-1cm] {Object student2};
\node (method) [draw, rectangle, fill=orange!20, below of=class, yshift=-4cm] {Call Method: greet / display};
\draw[->] (class) -- (object1);
\draw[->] (class) -- (object2);
\draw[->] (object1) -- (method);
\draw[->] (object2) -- (method);
\end{tikzpicture}
\end{center}

\section{Key Concepts}

\begin{itemize}
    \item Classes are \textbf{blueprints}; objects are \textbf{instances}  
    \item \texttt{\_\_init\_\_} initializes object attributes  
    \item Methods define actions objects can perform  
    \item Object attributes can be accessed and modified individually  
\end{itemize}

\section{End-of-Chapter Exercises}

\begin{exercise}{Very Short Questions}
\begin{enumerate}
    \item What is a class in Python?  
    \item What is an object?  
    \item What is the purpose of \texttt{\_\_init\_\_}?  
\end{enumerate}
\end{exercise}

\begin{exercise}{Short Answer Questions}
\begin{enumerate}
    \item Create a class \texttt{Car} with attributes \texttt{brand}, \texttt{model}, and \texttt{year}.  
    \item Write a method \texttt{info()} to print car details.  
    \item How do you create multiple objects of the same class?  
\end{enumerate}
\end{exercise}

\begin{exercise}{Activity-Based Questions}
\begin{enumerate}
    \item Create a class \texttt{Book} with attributes \texttt{title}, \texttt{author}, and \texttt{pages}. Make a method to display book info and create 3 book objects.  
    \item Modify the \texttt{Student} class to include a method \texttt{promote()} that increases grade by one. Test it for multiple students.  
    \item Create a program that stores employees with name, position, and salary. Loop through objects and display each employee’s details.  
\end{enumerate}
\end{exercise}

\section{Teacher Notes}

\marginnote{
\textbf{Learning Objectives:}  
\begin{itemize}
    \item Students understand classes, objects, attributes, and methods  
    \item Students can create multiple objects and access their data  
    \item Students can write simple methods and modify attributes  
\end{itemize}

\textbf{Teaching Tips:}  
\begin{itemize}
    \item Use real-life examples like students, cars, books, employees  
    \item Start with one class and one object, then expand to multiple objects  
    \item Demonstrate calling methods and modifying attributes live  
\end{itemize}

\textbf{Assessment:}  
\begin{itemize}
    \item Check correct creation of classes and objects  
    \item Observe understanding of methods and attribute manipulation  
\end{itemize}
}